\documentclass{amsart}
% For dealing with references we use the comment environment
\usepackage{verbatim}
\newenvironment{reference}{\comment}{\endcomment}
%\newenvironment{reference}{}{}
\newenvironment{slogan}{\comment}{\endcomment}
\newenvironment{history}{\comment}{\endcomment}

% For commutative diagrams we use Xy-pic
\usepackage[all]{xy}

% We use 2cell for 2-commutative diagrams.
\xyoption{2cell}
\UseAllTwocells

% We use multicol for the list of chapters between chapters
\usepackage{multicol}

% This is generally recommended for better output
\usepackage{lmodern}
\usepackage[T1]{fontenc}

% For cross-file-references

% Package for hypertext links:
\usepackage{hyperref}

% For any local file, say "hello.tex" you want to link to please

% Theorem environments.
%
\theoremstyle{plain}
\newtheorem{theorem}[subsection]{Theorem}
\newtheorem{proposition}[subsection]{Proposition}
\newtheorem{lemma}[subsection]{Lemma}

\theoremstyle{definition}
\newtheorem{definition}[subsection]{Definition}
\newtheorem{example}[subsection]{Example}
\newtheorem{exercise}[subsection]{Exercise}
\newtheorem{situation}[subsection]{Situation}

\theoremstyle{remark}
\newtheorem{remark}[subsection]{Remark}
\newtheorem{remarks}[subsection]{Remarks}

\numberwithin{equation}{subsection}

% Macros
%
\def\lim{\mathop{\mathrm{lim}}\nolimits}
\def\colim{\mathop{\mathrm{colim}}\nolimits}
\def\Spec{\mathop{\mathrm{Spec}}}
\def\Hom{\mathop{\mathrm{Hom}}\nolimits}
\def\Ext{\mathop{\mathrm{Ext}}\nolimits}
\def\SheafHom{\mathop{\mathcal{H}\!\mathit{om}}\nolimits}
\def\SheafExt{\mathop{\mathcal{E}\!\mathit{xt}}\nolimits}
\def\Sch{\mathit{Sch}}
\def\Mor{\mathop{\mathrm{Mor}}\nolimits}
\def\Ob{\mathop{\mathrm{Ob}}\nolimits}
\def\Sh{\mathop{\mathit{Sh}}\nolimits}
\def\NL{\mathop{N\!L}\nolimits}
\def\CH{\mathop{\mathrm{CH}}\nolimits}
\def\proetale{{pro\text{-}\acute{e}tale}}
\def\etale{{\acute{e}tale}}
\def\QCoh{\mathit{QCoh}}
\def\Ker{\mathop{\mathrm{Ker}}}
\def\Im{\mathop{\mathrm{Im}}}
\def\Coker{\mathop{\mathrm{Coker}}}
\def\Coim{\mathop{\mathrm{Coim}}}

% Boxtimes
%
\DeclareMathSymbol{\boxtimes}{\mathbin}{AMSa}{"02}

%
% Macros for moduli stacks/spaces
%
\def\QCohstack{\mathcal{QC}\!\mathit{oh}}
\def\Cohstack{\mathcal{C}\!\mathit{oh}}
\def\Spacesstack{\mathcal{S}\!\mathit{paces}}
\def\Quotfunctor{\mathrm{Quot}}
\def\Hilbfunctor{\mathrm{Hilb}}
\def\Curvesstack{\mathcal{C}\!\mathit{urves}}
\def\Polarizedstack{\mathcal{P}\!\mathit{olarized}}
\def\Complexesstack{\mathcal{C}\!\mathit{omplexes}}
% \Pic is the operator that assigns to X its picard group, usage \Pic(X)
% \Picardstack_{X/B} denotes the Picard stack of X over B
% \Picardfunctor_{X/B} denotes the Picard functor of X over B
\def\Pic{\mathop{\mathrm{Pic}}\nolimits}
\def\Picardstack{\mathcal{P}\!\mathit{ic}}
\def\Picardfunctor{\mathrm{Pic}}
\def\Deformationcategory{\mathcal{D}\!\mathit{ef}}

\setlength{\emergencystretch}{3em}
\begin{document}
\title[Stacks Project, Tag 07LG]{333 --- Crystalline cohomology of affine crystals is computed by the completed divided-power de Rham complex}
\maketitle
\noindent Copyright (C) 2005--2025 Johan de Jong. Stacks Project authors. Source: \href{https://stacks.math.columbia.edu/tag/07LG}{Tag 07LG}.

\noindent Editorial difficulty note (not author text). Difficulty H2: The proof must identify the crystal Cech resolution with a de Rham complex through a two-directional double-complex argument. It uses relative divided-power Poincare comparison for every row and verifies that all coprojection identifications agree on cohomology, producing the alternating zero/identity differentials and canonical comparison..

\noindent Editorial provenance note (not author text). History: excerpt from revision \texttt{a0444\allowbreak 6e57e\allowbreak c1fbc\allowbreak 252a8\allowbreak 71afc\allowbreak ec775\allowbreak 2fb28\allowbreak 07b14}, crystalline.tex, lines 3926--3988. Reference presentation was adapted; original-author context and core support blocks were retained or added. The mathematical wording is unchanged. Licensed under GNU FDL 1.2 or later; no invariant sections or cover texts. See ../licenses/COPYING. Context blocks reproduce author definitions and supporting results; they are not additional counted problems.

\begin{proposition}
\label{proposition-compute-cohomology}
With notations as above assume that
\begin{enumerate}
\item $\mathcal{F}$ is locally quasi-coherent, and
\item for any morphism $(U, T, \delta) \to (U', T', \delta')$
of $\text{Cris}(X/S)$ where $f : T \to T'$ is a closed immersion
the map $c_f : f^*\mathcal{F}_{T'} \to \mathcal{F}_T$ is surjective.
\end{enumerate}
Then the complex
$$
M(0) \to M(1) \to M(2) \to \ldots
$$
computes $R\Gamma(\text{Cris}(X/S), \mathcal{F})$.
\end{proposition}

\section{Cohomology in the affine case}
\label{section-cohomology-affine}

\noindent
Let's go back to the situation studied in
Section \ref{section-quasi-coherent-crystals}. We
start with $(A, I, \gamma)$ and $A/I \to C$ and set
$X = \Spec(C)$ and $S = \Spec(A)$. Then we choose
a polynomial ring $P$ over $A$ and a surjection $P \to C$ with
kernel $J$. We obtain $D$ and $D(n)$ see
(\ref{equation-D}) and (\ref{equation-Dn}).
Set $T(n)_e = \Spec(D(n)/p^eD(n))$ so that
$(X, T(n)_e, \delta(n))$ is an object of $\text{Cris}(X/S)$.
Let $\mathcal{F}$ be a sheaf of $\mathcal{O}_{X/S}$-modules and set
$$
M(n) = \lim_e \Gamma((X, T(n)_e, \delta(n)), \mathcal{F})
$$
for $n = 0, 1, 2, 3, \ldots$. This forms a cosimplicial module
over the cosimplicial ring $D(0), D(1), D(2), \ldots$.



\section{Crystals in quasi-coherent modules}
\label{section-quasi-coherent-crystals}

\noindent
In Situation \ref{situation-affine}.
Set $X = \Spec(C)$ and $S = \Spec(A)$. We are going to
classify crystals in quasi-coherent modules on $\text{Cris}(X/S)$.
Before we do so we fix some notation.

\medskip\noindent
Choose a polynomial ring $P = A[x_i]$ over $A$ and a surjection $P \to C$
of $A$-algebras with kernel $J = \Ker(P \to C)$. Set
\begin{equation}
\label{equation-D}
D = \lim_e D_{P, \gamma}(J) / p^eD_{P, \gamma}(J)
\end{equation}
for the $p$-adically completed divided power envelope.
This ring comes with a divided power ideal $\bar J$ and divided power
structure $\bar \gamma$, see Lemma \href{https://stacks.math.columbia.edu/tag/07KG}{lemma list properties (Tag 07KG)}.
Set $D_e = D/p^eD$ and denote $\bar J_e$ the image of $\bar J$ in $D_e$.
We will use the short hand
\begin{equation}
\label{equation-omega-D}
\Omega_D = \lim_e \Omega_{D_e/A, \bar\gamma} =
\lim_e \Omega_{D/A, \bar\gamma}/p^e\Omega_{D/A, \bar\gamma}
\end{equation}
for the $p$-adic completion of the module of divided power differentials,
see Lemma \href{https://stacks.math.columbia.edu/tag/07KK}{lemma differentials completion (Tag 07KK)}.
It is also the $p$-adic completion of
$\Omega_{D_{P, \gamma}(J)/A, \bar\gamma}$
which is free on $\text{d}x_i$, see
Lemma \href{https://stacks.math.columbia.edu/tag/07HW}{lemma module differentials divided power envelope (Tag 07HW)}.
Hence any element of $\Omega_D$ can be written uniquely as a sum
$\sum f_i\text{d}x_i$ with for all $e$ only finitely many $f_i$
not in $p^eD$. Moreover, the maps
$\text{d}_{D_e/A, \bar\gamma} : D_e \to \Omega_{D_e/A, \bar\gamma}$
fit together to define a divided power $A$-derivation
\begin{equation}
\label{equation-derivation-D}
\text{d} : D \longrightarrow \Omega_D
\end{equation}
on $p$-adic completions.

\medskip\noindent
We will also need the ``products $\Spec(D(n))$ of $\Spec(D)$'', see
Proposition \ref{proposition-compute-cohomology} and its proof for an
explanation. Formally these are defined as follows. For $n \geq 0$ let
$J(n) = \Ker(P \otimes_A \ldots \otimes_A P \to C)$ where
the tensor product has $n + 1$ factors. We set
\begin{equation}
\label{equation-Dn}
D(n) = \lim_e
D_{P \otimes_A \ldots \otimes_A P, \gamma}(J(n))/
p^eD_{P \otimes_A \ldots \otimes_A P, \gamma}(J(n))
\end{equation}
equal to the $p$-adic completion of the divided power envelope.
We denote $\bar J(n)$ its divided power ideal and $\bar \gamma(n)$
its divided powers. We also introduce $D(n)_e = D(n)/p^eD(n)$ as well
as the $p$-adically completed module of differentials
\begin{equation}
\label{equation-omega-Dn}
\Omega_{D(n)} = \lim_e \Omega_{D(n)_e/A, \bar\gamma} =
\lim_e \Omega_{D(n)/A, \bar\gamma}/p^e\Omega_{D(n)/A, \bar\gamma}
\end{equation}
and derivation
\begin{equation}
\label{equation-derivation-Dn}
\text{d} : D(n) \longrightarrow \Omega_{D(n)}
\end{equation}
Of course we have $D = D(0)$. Note that the rings $D(0), D(1), D(2), \ldots$
form a cosimplicial object in the category of divided power rings.



\noindent
In the lemma below we will consider pairs $(M, \nabla)$ satisfying the
following conditions
\begin{enumerate}
\item
\label{item-complete}
$M$ is a $p$-adically complete $D$-module,
\item
\label{item-connection}
$\nabla : M \to M \otimes^\wedge_D \Omega_D$ is a connection, i.e.,
$\nabla(fm) = m \otimes \text{d}f + f\nabla(m)$,
\item
\label{item-integrable}
$\nabla$ is integrable
(see Remark \href{https://stacks.math.columbia.edu/tag/07I0}{remark connection (Tag 07I0)}), and
\item
\label{item-topologically-quasi-nilpotent}
$\nabla$ is {\it topologically quasi-nilpotent}: If we write
$\nabla(m) = \sum \theta_i(m)\text{d}x_i$ for some operators
$\theta_i : M \to M$, then for any $m \in M$ there are only finitely
many pairs $(i, k)$ such that $\theta_i^k(m) \not \in pM$.
\end{enumerate}
The operators $\theta_i$ are sometimes denoted
$\nabla_{\partial/\partial x_i}$ in the literature.
In the following lemma we construct a functor from crystals in quasi-coherent
modules on $\text{Cris}(X/S)$ to the category of such pairs. We will show
this functor is an equivalence in
Proposition \href{https://stacks.math.columbia.edu/tag/07JH}{proposition crystals on affine (Tag 07JH)}.



\begin{situation}
\label{situation-affine}
Here $p$ is a prime number, $(A, I, \gamma)$ is a divided power
ring such that $A$ is a $\mathbf{Z}_{(p)}$-algebra, and $A \to C$ is a
ring map such that $IC = 0$ and such that $p$ is nilpotent in $C$.
\end{situation}

\begin{proposition}
\label{proposition-compute-cohomology-crystal}
Assumptions as in Proposition \ref{proposition-compute-cohomology}
but now assume that $\mathcal{F}$ is a crystal in quasi-coherent modules.
Let $(M, \nabla)$ be the corresponding module with connection over $D$, see
Proposition \href{https://stacks.math.columbia.edu/tag/07JH}{proposition crystals on affine (Tag 07JH)}. Then the complex
$$
M \otimes^\wedge_D \Omega^*_D
$$
computes $R\Gamma(\text{Cris}(X/S), \mathcal{F})$.
\end{proposition}

\begin{proof}
We will prove this using the two spectral sequences associated to the
double complex $K^{*, *}$ with terms
$$
K^{a, b} = M \otimes_D^\wedge \Omega^a_{D(b)}
$$
What do we know so far? Well, Lemma \href{https://stacks.math.columbia.edu/tag/07LA}{lemma vanishing (Tag 07LA)}
tells us that each column $K^{a, *}$, $a > 0$ is acyclic.
Proposition \ref{proposition-compute-cohomology} tells us that
the first column $K^{0, *}$ is quasi-isomorphic to
$R\Gamma(\text{Cris}(X/S), \mathcal{F})$.
Hence the first spectral sequence associated to the double complex
shows that there is a canonical quasi-isomorphism of
$R\Gamma(\text{Cris}(X/S), \mathcal{F})$ with
$\text{Tot}(K^{*, *})$.

\medskip\noindent
Next, let's consider the rows $K^{*, b}$. By
Lemma \href{https://stacks.math.columbia.edu/tag/07L2}{lemma structure Dn (Tag 07L2)}
each of the $b + 1$ maps $D \to D(b)$ presents $D(b)$ as the $p$-adic
completion of a divided power polynomial algebra over $D$.
Hence Lemma \href{https://stacks.math.columbia.edu/tag/07LD}{lemma relative poincare (Tag 07LD)} shows that the map
$$
M \otimes^\wedge_D\Omega^*_D
\longrightarrow
M \otimes^\wedge_{D(b)} \Omega^*_{D(b)} = K^{*, b}
$$
is a quasi-isomorphism. Note that each of these maps defines the {\it same}
map on cohomology (and even the same map in the derived category) as
the inverse is given by the co-diagonal map $D(b) \to D$ (corresponding
to the multiplication map $P \otimes_A \ldots \otimes_A P \to P$).
Hence if we look at the $E_1$ page of the second spectral sequence
we obtain
$$
E_1^{a, b} = H^a(M \otimes^\wedge_D\Omega^*_D)
$$
with differentials
$$
E_1^{a, 0} \xrightarrow{0}
E_1^{a, 1} \xrightarrow{1}
E_1^{a, 2} \xrightarrow{0}
E_1^{a, 3} \xrightarrow{1} \ldots
$$
as each of these is the alternation sum of the given identifications
$H^a(M \otimes^\wedge_D\Omega^*_D) = E_1^{a, 0} = E_1^{a, 1} = \ldots$.
Thus we see that the $E_2$ page is equal $H^a(M \otimes^\wedge_D\Omega^*_D)$
on the first row and zero elsewhere. It follows that the identification
of $M \otimes^\wedge_D\Omega^*_D$ with the first row induces a
quasi-isomorphism of $M \otimes^\wedge_D\Omega^*_D$ with
$\text{Tot}(K^{*, *})$.
\end{proof}
\end{document}
